
\begin{abstract}
Retrieving structured documents requires understanding both their content and how that content is organized. This study explores whether hyperbolic representations can help SANTA learn these relationships and improve dense retrieval. We focus on its Structured Data Alignment (SDA) component, which aligns structured data with natural-language text, and consider two hyperbolic representations: the Poincar\'e ball and the Lorentz hyperboloid. We plan to evaluate these variants on CodeSearchNet and Adv for code retrieval, alongside an S1000D XML retrieval dataset to be constructed for this study. Comparisons with the original SANTA model will use consistent encoder and training settings. The evaluation will examine retrieval effectiveness across the three datasets and hierarchical relationship preservation in S1000D XML. Our aim is to understand when geometric alignment helps retrieve structured technical documents and whether better retrieval is accompanied by better preservation of their structure.
\end{abstract}

\keywords{Dense retrieval; Hyperbolic representation learning; Geometric alignment; S1000D XML}
